\chapter{\MakeUppercase{Introdução}}
\label{cap:introducao}

\section{Contextualização}
% Apresentação do problema, exercícios resistidos, avaliação do movimento,
% sensores inerciais e motivação do HalterCheck.

\section{Motivação e justificativa}

\section{Problema e questão de pesquisa}

Este trabalho investiga a seguinte questão de pesquisa:

\begin{quote}
Em uma avaliação experimental intraindividual da elevação lateral bilateral, em que medida um sistema embarcado com IMUs em dois halteres consegue, em tempo quase real, distinguir execuções corretas de execuções com amplitude inadequada, alteração de cadência, assimetria entre os membros ou impulso/balanço, e quais erros posturais globais permanecem não identificáveis de forma confiável quando sensores corporais e imagens não são empregados como entradas do sistema?
\end{quote}

\section{Hipóteses}

\begin{description}
    \item[H1:] Sensores inerciais acoplados aos halteres permitem identificar
    amplitude inadequada na elevação lateral.

    \item[H2:] Sensores inerciais acoplados aos halteres permitem identificar
    alterações de cadência.

    \item[H3:] O uso de dois halteres instrumentados permite identificar
    assimetria entre membros.

    \item[H4:] Padrões de impulso ou balanço podem ser parcialmente
    identificados por picos e irregularidades no sinal inercial.

    \item[H5:] Erros posturais globais não são identificáveis de forma
    confiável apenas com IMUs nos halteres.
\end{description}

\section{Objetivos}

\subsection{Objetivo geral}

Avaliar a viabilidade de sensores inerciais acoplados a halteres para
identificar padrões corretos e erros específicos na execução da elevação
lateral, em tempo quase real, considerando as limitações de observabilidade
impostas pela ausência de sensores corporais ou câmeras.

\subsection{Objetivos específicos}

Para alcançar o objetivo geral, foram definidos os seguintes objetivos específicos:

\begin{enumerate}
    \item Desenvolver e caracterizar um protótipo composto por dois halteres instrumentados com IMUs, capazes de realizar aquisição sincronizada e transmissão dos sinais inerciais.
    
    \item Definir e executar um protocolo experimental intraindividual para a coleta de execuções da elevação lateral bilateral, contemplando o padrão de referência e variações controladas de amplitude, cadência, assimetria, impulso ou balanço e postura global.
    
    \item Construir um conjunto de dados rotulado e organizado por sessões independentes, contendo os sinais dos dois halteres, os metadados experimentais e as referências utilizadas na confirmação das condições executadas.
    
    \item Desenvolver um pipeline para calibração, sincronização, pré-processamento dos sinais e segmentação automática das repetições.
    
    \item Implementar métodos baseados em regras manuais e modelos de aprendizado de máquina como, por exemplo, Random Forest, Support Vector Machine e rede neural convolucional, para identificar os padrões de execução avaliados.
    
    \item Quantificar a capacidade dos métodos desenvolvidos de identificar alterações de amplitude, cadência, assimetria e impulso ou balanço, utilizando métricas de classificação e validação com separação entre sessões.
    
    \item Investigar os limites de observabilidade das IMUs acopladas exclusivamente aos halteres para a identificação de erros posturais globais, considerando uma referência externa de rotulagem.
    
    \item Comparar o desempenho, a estabilidade entre sessões e o custo computacional dos métodos baseados em regras e dos modelos de aprendizado de máquina.
    
    \item Mensurar a latência entre o término de uma repetição e a disponibilização do resultado, avaliando a viabilidade de geração de feedback em tempo quase real.

\end{enumerate}

A fim de assegurar que os objetivos específicos possam ser avaliados de
forma objetiva, cada um deles foi associado a evidências e métricas
observáveis. Essa associação estabelece a rastreabilidade entre os
objetivos, os procedimentos metodológicos e os resultados apresentados
posteriormente.

A matriz de rastreabilidade não pressupõe que os métodos desenvolvidos
alcançarão determinado desempenho mínimo. Sua finalidade é definir
previamente como cada objetivo será verificado, permitindo que resultados
positivos ou negativos contribuam para responder à questão de pesquisa.
O Quadro~\ref{qua:rastreabilidade-objetivos} apresenta essa relação.

\begin{quadro}[htb]
    \caption{Matriz de rastreabilidade dos objetivos específicos}
    \label{qua:rastreabilidade-objetivos}
    \centering
    \small
    \begin{tabular}{|p{0.12\textwidth}|p{0.78\textwidth}|}
        \hline
        \textbf{Objetivo} & \textbf{Evidências e métricas de verificação} \\
        \hline

        OE1 &
        Funcionamento dos dois halteres instrumentados; frequência efetiva
        de amostragem; perda de pacotes; erro residual de sincronização;
        estabilidade da comunicação e duração da operação contínua.
        \\
        \hline

        OE2 &
        Protocolo experimental documentado e executado; quantidade de
        sessões, séries e repetições válidas em cada condição experimental.
        \\
        \hline

        OE3 &
        Conjunto de dados organizado por sessões e condições; quantidade de
        repetições válidas; distribuição entre classes; presença dos
        metadados experimentais e dos respectivos rótulos.
        \\
        \hline

        OE4 &
        Funcionamento do pipeline de calibração, sincronização,
        pré-processamento e segmentação; precisão, revocação e medida F1
        na identificação das repetições; erro temporal na determinação do
        início e do término dos movimentos.
        \\
        \hline

        OE5 &
        Implementação e execução dos métodos baseados em regras manuais,
        Random Forest, Support Vector Machine e rede neural convolucional,
        utilizando as mesmas divisões de treinamento, validação e teste
        quando aplicável.
        \\
        \hline

        OE6 &
        Matrizes de confusão e métricas de precisão, revocação, medida F1
        e acurácia balanceada para amplitude, cadência, assimetria e
        impulso ou balanço, avaliadas em sessões não utilizadas no
        treinamento.
        \\
        \hline

        OE7 &
        Desempenho dos métodos na identificação dos erros posturais globais,
        comparação com uma referência externa de rotulagem e análise da
        estabilidade dos resultados entre sessões.
        \\
        \hline

        OE8 &
        Comparação entre os métodos quanto à medida F1 macro, acurácia
        balanceada, variabilidade entre sessões, tempo de inferência e
        recursos computacionais utilizados.
        \\
        \hline

        OE9 &
        Latência mediana, percentil 95, latência máxima e proporção de
        repetições cujo resultado foi disponibilizado dentro do limite
        definido para o feedback em tempo quase real.
        \\
        \hline
    \end{tabular}

    \begin{flushleft}
        \footnotesize Fonte: elaboração própria.
    \end{flushleft}
\end{quadro}

Essa matriz também orienta a organização das seções de Metodologia e
Resultados. Cada objetivo deverá possuir pelo menos um procedimento de
avaliação descrito na metodologia e uma evidência correspondente
apresentada nos resultados. Posteriormente, a discussão retomará essas
evidências para analisar o alcance dos objetivos e a sustentação das
hipóteses formuladas.

\newpage

\section{Escopo e delimitações}
\label{sec:escopo-delimitacoes}

Este trabalho está delimitado à avaliação experimental do HalterCheck
na elevação lateral bilateral com halteres. O estudo utiliza sessões
repetidas realizadas pelo próprio pesquisador, sob condições previamente
definidas de carga, postura inicial, amplitude e cadência. Os sinais
analisados são produzidos pelas IMUs acopladas aos dois halteres, sendo
cada repetição bilateral pareada adotada como unidade principal de
classificação.

Para manter o problema de pesquisa compatível com os recursos e o prazo
disponíveis, estão explicitamente fora do escopo deste trabalho:

\begin{itemize}
    \item realizar generalizações sobre o desempenho do sistema para
    outros indivíduos ou para uma população, uma vez que a avaliação é
    intraindividual;

    \item avaliar exercícios diferentes da elevação lateral bilateral ou
    modalidades executadas com barras, máquinas, cabos ou outros
    equipamentos;

    \item avaliar os efeitos do exercício sobre força, hipertrofia,
    condicionamento físico, dor ou risco de lesão;

    \item realizar diagnóstico clínico, prescrição de treinamento ou
    substituir a avaliação de profissionais de educação física,
    fisioterapia ou medicina;

    \item reconstruir integralmente a postura corporal ou estimar todos
    os ângulos articulares do participante;

    \item utilizar câmeras ou sensores corporais como entradas dos
    métodos de identificação desenvolvidos. Uma referência externa
    poderá ser utilizada exclusivamente para confirmar e rotular as
    condições executadas;

    \item estimar diretamente e com precisão clínica o ângulo de
    abdução do ombro. Os ângulos definidos no protocolo serão utilizados
    como referências operacionais para a produção e rotulagem das
    condições experimentais;

    \item avaliar todas as combinações possíveis entre erros. As
    condições principais serão executadas isoladamente para evitar
    classes combinadas com quantidade insuficiente de observações;

    \item investigar de forma sistemática os efeitos de fadiga, variação
    de carga, experiência de treinamento ou alterações anatômicas sobre
    os sinais inerciais;

    \item desenvolver modelos generalizáveis ou personalizados para
    diferentes usuários, uma vez que não haverá dados de outros
    participantes;

    \item fornecer correções contínuas durante a execução da repetição.
    O feedback em tempo quase real será considerado como o resultado
    disponibilizado após a conclusão de cada repetição;

    \item exigir que todo o processamento e a inferência sejam
    executados nas ESP32. O processamento poderá ser realizado por um
    servidor conectado aos dispositivos;

    \item desenvolver um produto comercial final, incluindo aplicativo
    para usuários, certificação, fabricação em escala, impermeabilização
    ou validação de durabilidade em longo prazo.
\end{itemize}

Essas delimitações não impedem que alguns dos aspectos mencionados sejam
explorados futuramente. Elas estabelecem os limites dentro dos quais os
resultados e as conclusões desta pesquisa deverão ser interpretados.

\section{Contribuições do trabalho}
\label{sec:contribuicoes}

As contribuições deste trabalho são organizadas em duas categorias. A
contribuição científica está relacionada ao conhecimento produzido pela
avaliação experimental, enquanto a contribuição de engenharia está
relacionada ao artefato desenvolvido e às soluções técnicas necessárias
para sua operação.

\subsection{Contribuição científica}
\label{subsec:contribuicao-cientifica}

A contribuição científica principal consiste na produção de evidências
experimentais sobre a capacidade e os limites de sensores inerciais
acoplados exclusivamente a dois halteres para identificar características
da elevação lateral bilateral.

O estudo busca caracterizar quais informações relacionadas à amplitude,
cadência, assimetria e uso de impulso ou balanço permanecem observáveis
nos sinais dos halteres e quais erros posturais globais não podem ser
identificados de forma estável na configuração avaliada. Essa análise
contribui para delimitar o potencial de sensores posicionados no
equipamento, em contraste com abordagens que dependem de câmeras ou
sensores fixados ao corpo.

Também constitui contribuição científica a comparação experimental
entre métodos baseados em regras manuais, Random Forest, Support Vector
Machine e rede neural convolucional, utilizando separação entre sessões
para avaliar a estabilidade dos resultados em dados não empregados no
treinamento.

Como contribuição metodológica complementar, o trabalho estabelece um
protocolo intraindividual para produção, rotulagem e avaliação de
execuções controladas da elevação lateral bilateral com dois halteres
instrumentados.

\subsection{Contribuição de engenharia}
\label{subsec:contribuicao-engenharia}

A contribuição de engenharia consiste no desenvolvimento do HalterCheck,
um protótipo formado por dois halteres instrumentados com IMUs e
unidades de processamento independentes, integrado a uma infraestrutura
para aquisição bilateral, sincronização temporal e transmissão dos
sinais inerciais.

O protótipo inclui firmware para leitura e identificação dos sensores,
mecanismos para comunicação dos dispositivos com o servidor e um
pipeline para controle de qualidade, pré-processamento, segmentação,
contagem e classificação das repetições. A solução também permite medir
a latência entre o término do movimento e a produção do resultado,
possibilitando avaliar sua aplicação em feedback próximo do tempo real.

A documentação da arquitetura, dos formatos de dados, dos procedimentos
experimentais e dos métodos de processamento busca tornar o sistema
reproduzível e permitir sua extensão em investigações futuras.


\section{Organização do trabalho}
% Resumo dos capítulos seguintes.